
\documentclass{article}
\usepackage{colortbl}
\usepackage{makecell}
\usepackage{multirow}
\usepackage{supertabular}

\begin{document}

\newcounter{utterance}

\centering \large Interaction Transcript for game `codenames', experiment `opponent\_difficult', episode 4 with nemotron{-}3{-}ultra.
\vspace{24pt}

{ \footnotesize  \setcounter{utterance}{1}
\setlength{\tabcolsep}{0pt}
\begin{supertabular}{c@{$\;$}|p{.15\linewidth}@{}p{.15\linewidth}p{.15\linewidth}p{.15\linewidth}p{.15\linewidth}p{.15\linewidth}}
    \# & $\;$A & \multicolumn{4}{c}{Game Master} & $\;\:$B\\
    \hline

    \theutterance \stepcounter{utterance}  
    & & & \multicolumn{2}{p{0.3\linewidth}}{
        \cellcolor[rgb]{0.9,0.9,0.9}{
            \makecell[{{p{\linewidth}}}]{
                \texttt{\tiny{[GM$|$GM]}}
                {'hidden': {'team': ['قش', 'أوروبا', 'عقد', 'وردة', 'سياج', 'موجة', 'زلة', 'ديناصور', 'كازينو'], 'innocent': ['قدم', 'طائرة نفاثة', 'نيويورك', 'شاطئ', 'سند', 'دائرة', 'الفندق'], 'opponent': ['الأمازون', 'زر', 'زجاجة', 'ليمون', 'مؤامرة', 'تاج', 'فطيرة', 'طوكيو'], 'assassin': ['رأس المال']}, 'revealed': {'team': {'team': [], 'innocent': [], 'opponent': [], 'assassin': []}, 'opponent': {'team': [], 'innocent': [], 'opponent': [], 'assassin': []}}}
            }
        }
    }
    & & \\ \\

    \theutterance \stepcounter{utterance}  
    & & \multicolumn{4}{p{0.6\linewidth}}{
        \cellcolor[rgb]{0.9,0.9,0.9}{
            \makecell[{{p{\linewidth}}}]{
                \texttt{\tiny{[P1$\langle$GM]}}
                \texttt{قدّم تلميحًا واحدًا مكوّنًا من كلمة واحدة يرتبط ببعض كلمات فريقك من القائمة أدناه. يمكنك اختيار عدد الكلمات التي يرتبط بها تلميحك وأيّها.} \\
\texttt{الهدف هو أن يخمّن شريكك جميع كلمات فريقك بناءً على التلميحات التي تقدّمها له، قبل أن يكشف الفريق الآخر جميع كلماته. إذا خَمَّن شريكك كلمة 'الخصم' أو كلمة 'المشتّت'، تنتهي جولة التخمين الخاصة به ويحصل الفريق المنافس على دور. إذا خَمَّن شريكك كلمة 'القاتل'، تنتهي اللعبة ويخسر فريقك فورًا.} \\
\texttt{كلما ارتبط تلميحك بعدد أكبر من الكلمات، كان بإمكانك الفوز أسرع، لكن قد تصبح الروابط أصعب على شريكك في التخمين، لذا اختر كلمة التلميح بحكمة. يجب أن تكون كلمة التلميح مرتبطة دلاليًا بالكلمات المستهدفة، ولا يمكن أن تكون واحدة من الكلمات في القوائم أو أن تحتوي على أجزاء منها.} \\
\\ 
\texttt{قدّم دائمًا تلميحك المكوّن من كلمة واحدة وقائمتك المفصولة بفواصل من الكلمات المستهدفة المرتبطة بالتنسيق التالي واجعل إجاباتك قصيرة قدر الإمكان، ولا تُدرج أي نص آخر غير المطلوب في هذا النموذج:} \\
\\ 
\texttt{تلميح: <كلمة>} \\
\texttt{أهداف: <أهداف>} \\
\\ 
\texttt{كلمات فريقك هي: قش, أوروبا, عقد, وردة, سياج, موجة, زلة, ديناصور, كازينو.} \\
\texttt{كلمات فريق خصمك هي: الأمازون, زر, زجاجة, ليمون, مؤامرة, تاج, فطيرة, طوكيو.} \\
\texttt{كلمات التشتيت هي: قدم, طائرة نفاثة, نيويورك, شاطئ, سند, دائرة, الفندق.} \\
\texttt{كلمات القتل هي: رأس المال.} \\
            }
        }
    }
    & \\ \\

    \theutterance \stepcounter{utterance}  
    & \multicolumn{4}{p{0.6\linewidth}}{
        \cellcolor[rgb]{0.9,0.9,0.9}{
            \makecell[{{p{\linewidth}}}]{
                \texttt{\tiny{[P1$\rangle$GM]}}
                \texttt{تلميح: عصر} \\
\texttt{أهداف: أوروبا, عقد, ديناصور} \\
            }
        }
    }
    & & \\ \\

    \theutterance \stepcounter{utterance}  
    & & & \multicolumn{2}{p{0.3\linewidth}}{
        \cellcolor[rgb]{0.9,0.9,0.9}{
            \makecell[{{p{\linewidth}}}]{
                \texttt{\tiny{[GM$|$GM]}}
                \texttt{عصر} \\
            }
        }
    }
    & & \\ \\

    \theutterance \stepcounter{utterance}  
    & & & \multicolumn{2}{p{0.3\linewidth}}{
        \cellcolor[rgb]{0.9,0.9,0.9}{
            \makecell[{{p{\linewidth}}}]{
                \texttt{\tiny{[GM$|$GM]}}
                ['أوروبا', 'عقد', 'ديناصور']
            }
        }
    }
    & & \\ \\

    \theutterance \stepcounter{utterance}  
    & & & \multicolumn{2}{p{0.3\linewidth}}{
        \cellcolor[rgb]{0.9,0.9,0.9}{
            \makecell[{{p{\linewidth}}}]{
                \texttt{\tiny{[GM$|$GM]}}
                {'word': 'أوروبا', 'assignment': 'team'}
            }
        }
    }
    & & \\ \\

    \theutterance \stepcounter{utterance}  
    & & & \multicolumn{2}{p{0.3\linewidth}}{
        \cellcolor[rgb]{0.9,0.9,0.9}{
            \makecell[{{p{\linewidth}}}]{
                \texttt{\tiny{[GM$|$GM]}}
                {'word': 'عقد', 'assignment': 'team'}
            }
        }
    }
    & & \\ \\

    \theutterance \stepcounter{utterance}  
    & & & \multicolumn{2}{p{0.3\linewidth}}{
        \cellcolor[rgb]{0.9,0.9,0.9}{
            \makecell[{{p{\linewidth}}}]{
                \texttt{\tiny{[GM$|$GM]}}
                {'word': 'ديناصور', 'assignment': 'team'}
            }
        }
    }
    & & \\ \\

    \theutterance \stepcounter{utterance}  
    & & \multicolumn{4}{p{0.6\linewidth}}{
        \cellcolor[rgb]{0.9,0.9,0.9}{
            \makecell[{{p{\linewidth}}}]{
                \texttt{\tiny{[GM$\rangle$P2]}}
                \texttt{قدّم قائمة مفصولة بفواصل تضم ما يصل إلى 3 من الكلمات من القائمة التالية التي ترتبط على أفضل وجه أو تكون الأكثر ارتباطًا بالكلمة 'عصر'. ابدأ دائمًا قائمة تخميناتك بـ تخمين:  ولا تُدرج أي نص آخر في إجابتك.} \\
\\ 
\texttt{أوروبا, رأس المال, سياج, كازينو, قدم, الأمازون, سند, زجاجة, قش, طوكيو, زر, دائرة, ليمون, مؤامرة, نيويورك, شاطئ, موجة, الفندق, وردة, زلة, طائرة نفاثة, تاج, عقد, ديناصور, فطيرة} \\
            }
        }
    }
    & \\ \\

    \theutterance \stepcounter{utterance}  
    & & & \multicolumn{4}{p{0.6\linewidth}}{
        \cellcolor[rgb]{0.9,0.9,0.9}{
            \makecell[{{p{\linewidth}}}]{
                \texttt{\tiny{[GM$\langle$P2]}}
                \texttt{تخمين: ديناصور, عقد, فطيرة} \\
            }
        }
    }
     \\ \\

    \theutterance \stepcounter{utterance}  
    & & & \multicolumn{2}{p{0.3\linewidth}}{
        \cellcolor[rgb]{0.9,0.9,0.9}{
            \makecell[{{p{\linewidth}}}]{
                \texttt{\tiny{[GM$|$GM]}}
                ['ديناصور', 'عقد', 'فطيرة']
            }
        }
    }
    & & \\ \\

    \theutterance \stepcounter{utterance}  
    & & & \multicolumn{2}{p{0.3\linewidth}}{
        \cellcolor[rgb]{0.9,0.9,0.9}{
            \makecell[{{p{\linewidth}}}]{
                \texttt{\tiny{[GM$|$GM]}}
                {'word': 'ديناصور', 'assignment': 'team'}
            }
        }
    }
    & & \\ \\

    \theutterance \stepcounter{utterance}  
    & & & \multicolumn{2}{p{0.3\linewidth}}{
        \cellcolor[rgb]{0.9,0.9,0.9}{
            \makecell[{{p{\linewidth}}}]{
                \texttt{\tiny{[GM$|$GM]}}
                {'word': 'ديناصور', 'assignment': 'team'}
            }
        }
    }
    & & \\ \\

    \theutterance \stepcounter{utterance}  
    & & & \multicolumn{2}{p{0.3\linewidth}}{
        \cellcolor[rgb]{0.9,0.9,0.9}{
            \makecell[{{p{\linewidth}}}]{
                \texttt{\tiny{[GM$|$GM]}}
                {'word': 'عقد', 'assignment': 'team'}
            }
        }
    }
    & & \\ \\

    \theutterance \stepcounter{utterance}  
    & & & \multicolumn{2}{p{0.3\linewidth}}{
        \cellcolor[rgb]{0.9,0.9,0.9}{
            \makecell[{{p{\linewidth}}}]{
                \texttt{\tiny{[GM$|$GM]}}
                {'word': 'عقد', 'assignment': 'team'}
            }
        }
    }
    & & \\ \\

    \theutterance \stepcounter{utterance}  
    & & & \multicolumn{2}{p{0.3\linewidth}}{
        \cellcolor[rgb]{0.9,0.9,0.9}{
            \makecell[{{p{\linewidth}}}]{
                \texttt{\tiny{[GM$|$GM]}}
                {'word': 'فطيرة', 'assignment': 'opponent'}
            }
        }
    }
    & & \\ \\

    \theutterance \stepcounter{utterance}  
    & & & \multicolumn{2}{p{0.3\linewidth}}{
        \cellcolor[rgb]{0.9,0.9,0.9}{
            \makecell[{{p{\linewidth}}}]{
                \texttt{\tiny{[GM$|$GM]}}
                \texttt{فطيرة} \\
            }
        }
    }
    & & \\ \\

    \theutterance \stepcounter{utterance}  
    & & & \multicolumn{2}{p{0.3\linewidth}}{
        \cellcolor[rgb]{0.9,0.9,0.9}{
            \makecell[{{p{\linewidth}}}]{
                \texttt{\tiny{[GM$|$GM]}}
                {'word': 'تاج', 'assignment': 'opponent'}
            }
        }
    }
    & & \\ \\

    \theutterance \stepcounter{utterance}  
    & & & \multicolumn{2}{p{0.3\linewidth}}{
        \cellcolor[rgb]{0.9,0.9,0.9}{
            \makecell[{{p{\linewidth}}}]{
                \texttt{\tiny{[GM$|$GM]}}
                {'word': 'مؤامرة', 'assignment': 'opponent'}
            }
        }
    }
    & & \\ \\

    \theutterance \stepcounter{utterance}  
    & & & \multicolumn{2}{p{0.3\linewidth}}{
        \cellcolor[rgb]{0.9,0.9,0.9}{
            \makecell[{{p{\linewidth}}}]{
                \texttt{\tiny{[GM$|$GM]}}
                {'hidden': {'team': ['قش', 'أوروبا', 'وردة', 'سياج', 'موجة', 'زلة', 'كازينو'], 'innocent': ['قدم', 'طائرة نفاثة', 'نيويورك', 'شاطئ', 'سند', 'دائرة', 'الفندق'], 'opponent': ['الأمازون', 'زر', 'زجاجة', 'ليمون', 'طوكيو'], 'assassin': ['رأس المال']}, 'revealed': {'team': {'team': ['ديناصور', 'عقد'], 'innocent': [], 'opponent': ['فطيرة'], 'assassin': []}, 'opponent': {'team': [], 'innocent': [], 'opponent': ['تاج', 'مؤامرة'], 'assassin': []}}}
            }
        }
    }
    & & \\ \\

    \theutterance \stepcounter{utterance}  
    & & \multicolumn{4}{p{0.6\linewidth}}{
        \cellcolor[rgb]{0.9,0.9,0.9}{
            \makecell[{{p{\linewidth}}}]{
                \texttt{\tiny{[P1$\langle$GM]}}
                \texttt{تم تخمين الكلمات ديناصور, عقد بشكل صحيح. تم تخمين الكلمة فطيرة، لكنها كلمة الخصم. انتهى دور زميلك في الفريق هنا.} \\
\texttt{قدّم الآن تلميحًا آخر يرتبط ببعض كلمات فريقك المتبقية وقائمة بالكلمات الهدف ذات الصلة. تذكّر أن تبدأ تلميحك بـ 'تلميح: '، ثم ضع سطرًا جديدًا، وابدأ قائمتك المفصولة بفواصل للكلمات الهدف بـ 'أهداف: '. ملاحظة: تمت إزالة بعض الكلمات من القوائم مقارنة بالطلبات السابقة.} \\
\\ 
\texttt{كلمات فريقك المتبقية هي: قش, أوروبا, وردة, سياج, موجة, زلة, كازينو.} \\
\texttt{الكلمات المتبقية لخصمك هي: الأمازون, زر, زجاجة, ليمون, مؤامرة, تاج, طوكيو.} \\
\texttt{كلمات التشتيت المتبقية هي: قدم, طائرة نفاثة, نيويورك, شاطئ, سند, دائرة, الفندق.} \\
\texttt{كلمات القتل المتبقية هي: رأس المال.} \\
            }
        }
    }
    & \\ \\

    \theutterance \stepcounter{utterance}  
    & \multicolumn{4}{p{0.6\linewidth}}{
        \cellcolor[rgb]{0.9,0.9,0.9}{
            \makecell[{{p{\linewidth}}}]{
                \texttt{\tiny{[P1$\rangle$GM]}}
                \texttt{تلميح: نبتة} \\
\texttt{أهداف: قش, وردة, سياج} \\
            }
        }
    }
    & & \\ \\

    \theutterance \stepcounter{utterance}  
    & & & \multicolumn{2}{p{0.3\linewidth}}{
        \cellcolor[rgb]{0.9,0.9,0.9}{
            \makecell[{{p{\linewidth}}}]{
                \texttt{\tiny{[GM$|$GM]}}
                \texttt{نبتة} \\
            }
        }
    }
    & & \\ \\

    \theutterance \stepcounter{utterance}  
    & & & \multicolumn{2}{p{0.3\linewidth}}{
        \cellcolor[rgb]{0.9,0.9,0.9}{
            \makecell[{{p{\linewidth}}}]{
                \texttt{\tiny{[GM$|$GM]}}
                ['قش', 'وردة', 'سياج']
            }
        }
    }
    & & \\ \\

    \theutterance \stepcounter{utterance}  
    & & & \multicolumn{2}{p{0.3\linewidth}}{
        \cellcolor[rgb]{0.9,0.9,0.9}{
            \makecell[{{p{\linewidth}}}]{
                \texttt{\tiny{[GM$|$GM]}}
                {'word': 'قش', 'assignment': 'team'}
            }
        }
    }
    & & \\ \\

    \theutterance \stepcounter{utterance}  
    & & & \multicolumn{2}{p{0.3\linewidth}}{
        \cellcolor[rgb]{0.9,0.9,0.9}{
            \makecell[{{p{\linewidth}}}]{
                \texttt{\tiny{[GM$|$GM]}}
                {'word': 'وردة', 'assignment': 'team'}
            }
        }
    }
    & & \\ \\

    \theutterance \stepcounter{utterance}  
    & & & \multicolumn{2}{p{0.3\linewidth}}{
        \cellcolor[rgb]{0.9,0.9,0.9}{
            \makecell[{{p{\linewidth}}}]{
                \texttt{\tiny{[GM$|$GM]}}
                {'word': 'سياج', 'assignment': 'team'}
            }
        }
    }
    & & \\ \\

    \theutterance \stepcounter{utterance}  
    & & \multicolumn{4}{p{0.6\linewidth}}{
        \cellcolor[rgb]{0.9,0.9,0.9}{
            \makecell[{{p{\linewidth}}}]{
                \texttt{\tiny{[GM$\rangle$P2]}}
                \texttt{تم تخمين الكلمات ديناصور, عقد بشكل صحيح. تم تخمين الكلمة فطيرة، لكنها كلمة الخصم. انتهى دورك هنا.} \\
\texttt{قدّم الآن قائمة أخرى مفصولة بفواصل تتضمن ما لا يقل عن 1 وما يصل إلى 3 من الكلمات من قائمة الكلمات التالية التي ترتبط على أفضل وجه أو تكون الأكثر ارتباطًا بالكلمة 'نبتة'. تذكّر أن تبدأ إجابتك بـ تخمين: . ملاحظة: تمت إزالة بعض الكلمات من القائمة مقارنة بالطلبات السابقة.} \\
\\ 
\texttt{أوروبا, رأس المال, سياج, كازينو, قدم, الأمازون, سند, زجاجة, قش, طوكيو, زر, دائرة, ليمون, نيويورك, شاطئ, موجة, الفندق, وردة, زلة, طائرة نفاثة} \\
            }
        }
    }
    & \\ \\

    \theutterance \stepcounter{utterance}  
    & & & \multicolumn{4}{p{0.6\linewidth}}{
        \cellcolor[rgb]{0.9,0.9,0.9}{
            \makecell[{{p{\linewidth}}}]{
                \texttt{\tiny{[GM$\langle$P2]}}
                \texttt{تخمين: وردة, قش, الأمازون} \\
            }
        }
    }
     \\ \\

    \theutterance \stepcounter{utterance}  
    & & & \multicolumn{2}{p{0.3\linewidth}}{
        \cellcolor[rgb]{0.9,0.9,0.9}{
            \makecell[{{p{\linewidth}}}]{
                \texttt{\tiny{[GM$|$GM]}}
                ['وردة', 'قش', 'الأمازون']
            }
        }
    }
    & & \\ \\

    \theutterance \stepcounter{utterance}  
    & & & \multicolumn{2}{p{0.3\linewidth}}{
        \cellcolor[rgb]{0.9,0.9,0.9}{
            \makecell[{{p{\linewidth}}}]{
                \texttt{\tiny{[GM$|$GM]}}
                {'word': 'وردة', 'assignment': 'team'}
            }
        }
    }
    & & \\ \\

    \theutterance \stepcounter{utterance}  
    & & & \multicolumn{2}{p{0.3\linewidth}}{
        \cellcolor[rgb]{0.9,0.9,0.9}{
            \makecell[{{p{\linewidth}}}]{
                \texttt{\tiny{[GM$|$GM]}}
                {'word': 'وردة', 'assignment': 'team'}
            }
        }
    }
    & & \\ \\

    \theutterance \stepcounter{utterance}  
    & & & \multicolumn{2}{p{0.3\linewidth}}{
        \cellcolor[rgb]{0.9,0.9,0.9}{
            \makecell[{{p{\linewidth}}}]{
                \texttt{\tiny{[GM$|$GM]}}
                {'word': 'قش', 'assignment': 'team'}
            }
        }
    }
    & & \\ \\

    \theutterance \stepcounter{utterance}  
    & & & \multicolumn{2}{p{0.3\linewidth}}{
        \cellcolor[rgb]{0.9,0.9,0.9}{
            \makecell[{{p{\linewidth}}}]{
                \texttt{\tiny{[GM$|$GM]}}
                {'word': 'قش', 'assignment': 'team'}
            }
        }
    }
    & & \\ \\

    \theutterance \stepcounter{utterance}  
    & & & \multicolumn{2}{p{0.3\linewidth}}{
        \cellcolor[rgb]{0.9,0.9,0.9}{
            \makecell[{{p{\linewidth}}}]{
                \texttt{\tiny{[GM$|$GM]}}
                {'word': 'الأمازون', 'assignment': 'opponent'}
            }
        }
    }
    & & \\ \\

    \theutterance \stepcounter{utterance}  
    & & & \multicolumn{2}{p{0.3\linewidth}}{
        \cellcolor[rgb]{0.9,0.9,0.9}{
            \makecell[{{p{\linewidth}}}]{
                \texttt{\tiny{[GM$|$GM]}}
                \texttt{الأمازون} \\
            }
        }
    }
    & & \\ \\

    \theutterance \stepcounter{utterance}  
    & & & \multicolumn{2}{p{0.3\linewidth}}{
        \cellcolor[rgb]{0.9,0.9,0.9}{
            \makecell[{{p{\linewidth}}}]{
                \texttt{\tiny{[GM$|$GM]}}
                {'word': 'ليمون', 'assignment': 'opponent'}
            }
        }
    }
    & & \\ \\

    \theutterance \stepcounter{utterance}  
    & & & \multicolumn{2}{p{0.3\linewidth}}{
        \cellcolor[rgb]{0.9,0.9,0.9}{
            \makecell[{{p{\linewidth}}}]{
                \texttt{\tiny{[GM$|$GM]}}
                {'word': 'طوكيو', 'assignment': 'opponent'}
            }
        }
    }
    & & \\ \\

    \theutterance \stepcounter{utterance}  
    & & & \multicolumn{2}{p{0.3\linewidth}}{
        \cellcolor[rgb]{0.9,0.9,0.9}{
            \makecell[{{p{\linewidth}}}]{
                \texttt{\tiny{[GM$|$GM]}}
                {'hidden': {'team': ['أوروبا', 'سياج', 'موجة', 'زلة', 'كازينو'], 'innocent': ['قدم', 'طائرة نفاثة', 'نيويورك', 'شاطئ', 'سند', 'دائرة', 'الفندق'], 'opponent': ['زر', 'زجاجة'], 'assassin': ['رأس المال']}, 'revealed': {'team': {'team': ['ديناصور', 'عقد', 'وردة', 'قش'], 'innocent': [], 'opponent': ['فطيرة', 'الأمازون'], 'assassin': []}, 'opponent': {'team': [], 'innocent': [], 'opponent': ['تاج', 'مؤامرة', 'ليمون', 'طوكيو'], 'assassin': []}}}
            }
        }
    }
    & & \\ \\

    \theutterance \stepcounter{utterance}  
    & & \multicolumn{4}{p{0.6\linewidth}}{
        \cellcolor[rgb]{0.9,0.9,0.9}{
            \makecell[{{p{\linewidth}}}]{
                \texttt{\tiny{[P1$\langle$GM]}}
                \texttt{تم تخمين الكلمات وردة, قش بشكل صحيح. تم تخمين الكلمة الأمازون، لكنها كلمة الخصم. انتهى دور زميلك في الفريق هنا.} \\
\texttt{قدّم الآن تلميحًا آخر يرتبط ببعض كلمات فريقك المتبقية وقائمة بالكلمات الهدف ذات الصلة. تذكّر أن تبدأ تلميحك بـ 'تلميح: '، ثم ضع سطرًا جديدًا، وابدأ قائمتك المفصولة بفواصل للكلمات الهدف بـ 'أهداف: '. ملاحظة: تمت إزالة بعض الكلمات من القوائم مقارنة بالطلبات السابقة.} \\
\\ 
\texttt{كلمات فريقك المتبقية هي: أوروبا, سياج, موجة, زلة, كازينو.} \\
\texttt{الكلمات المتبقية لخصمك هي: زر, زجاجة, ليمون, طوكيو.} \\
\texttt{كلمات التشتيت المتبقية هي: قدم, طائرة نفاثة, نيويورك, شاطئ, سند, دائرة, الفندق.} \\
\texttt{كلمات القتل المتبقية هي: رأس المال.} \\
            }
        }
    }
    & \\ \\

    \theutterance \stepcounter{utterance}  
    & \multicolumn{4}{p{0.6\linewidth}}{
        \cellcolor[rgb]{0.9,0.9,0.9}{
            \makecell[{{p{\linewidth}}}]{
                \texttt{\tiny{[P1$\rangle$GM]}}
                \texttt{تلميح: سياج} \\
\texttt{أهداف: سياج} \\
            }
        }
    }
    & & \\ \\

    \theutterance \stepcounter{utterance}  
    & & & \multicolumn{2}{p{0.3\linewidth}}{
        \cellcolor[rgb]{0.9,0.9,0.9}{
            \makecell[{{p{\linewidth}}}]{
                \texttt{\tiny{[GM$|$GM]}}
                {'player': 'Cluegiver', 'type': 'clue is morphologically related to word on the board', 'utterance': 'تلميح: سياج\nأهداف: سياج', 'message': "Your clue 'سياج' is morphologically similar to the word سياج, please choose another clue word.", 'clue': 'سياج', 'similar_board_word': 'سياج'}
            }
        }
    }
    & & \\ \\

\end{supertabular}
}

\end{document}
